\section{Domande di ricerca}

Le domande di ricerca sono mantenute in numero ridotto e seguono la distinzione tra presenza del segnale, eterogeneita' per esposizione, capacita' predittiva e segnali alternativi.

\begin{description}
  \item[RQ1 -- Presenza e temporalita' del segnale.] Gli account successivamente sottoposti a takedown ricevono piu' blocchi dei controlli comparabili? In quale momento della finestra pre-evento emerge la differenza?
  \item[RQ2 -- Eterogeneita' per esposizione.] Il potere informativo dei blocchi cambia in funzione dell'attivita' e della visibilita' dell'account, con particolare attenzione ai gruppi 00, 0p, p0 e pp?
  \item[RQ3 -- Potere predittivo.] Quanto le feature relative ai blocchi ricevuti consentono a una Random Forest di distinguere gli account positivi dai controlli negativi?
  \item[RQ4 -- Segnali alternativi.] Labels ufficiali, labeler di terze parti e modlist hanno coverage e capacita' discriminante sufficienti per un utilizzo predittivo sistematico?
\end{description}

L'affidabilita' della label \code{!takedown} e' trattata come controllo preliminare sulla validita' del target. Non e' formulata come domanda autonoma perche' serve soprattutto a stabilire se il target operativo sia abbastanza stabile da supportare l'analisi successiva.
